\let\negmedspace\undefined
\let\negthickspace\undefined

\documentclass[journal]{IEEEtran}
\usepackage[a5paper, margin=10mm, onecolumn]{geometry}
\usepackage{tfrupee}

\setlength{\headheight}{1cm}
\setlength{\headsep}{0mm}

\usepackage{gvv-book}
\usepackage{cite}
\usepackage{amsmath,amssymb,amsfonts,amsthm}
\usepackage{algorithmic}
\usepackage{graphicx}
\usepackage{textcomp}
\usepackage{xcolor}
\usepackage{txfonts}
\usepackage{listings}
\usepackage{enumitem}
\usepackage{mathtools}
\usepackage{gensymb}
\usepackage{comment}
\usepackage[breaklinks=true]{hyperref}
\usepackage{tkz-euclide}
\usepackage{listings}
\usepackage{gvv}
\usepackage{tikz}
\def\inputGnumericTable{}
\usepackage[latin1]{inputenc}
\usepackage{color}
\usepackage{array}
\usepackage{longtable}
\usepackage{calc}
\usepackage{multirow}
\usepackage{hhline}
\usepackage{ifthen}
\usepackage{lscape}

\begin{document}

\bibliographystyle{IEEEtran}

\title{2.4.6}
\author{EE25BTECH11016 - Taraka Abhinav}
{\let\newpage\relax\maketitle}

\renewcommand{\thefigure}{\theenumi}
\renewcommand{\thetable}{\theenumi}
\setlength{\intextsep}{10pt}

\numberwithin{figure}{enumi}
\renewcommand{\thetable}{\theenumi}

\textbf{Question}:\\
If $\vec{a} = 5\hat{i} - \hat{j} - 3\hat{k}$ and $\vec{b} = \hat{i} + 3\hat{j} - 5\hat{k}$, then show that the vectors $\vec{a} + \vec{b}$ and $\vec{a} - \vec{b}$ are perpendicular. \\

\solution \\

Two vectors are perpendicular if their dot product is zero. We need to show that $(\vec{a} + \vec{b})^T (\vec{a} - \vec{b}) = 0$.

First, we expand the algebraic expression for the dot product:
\begin{align}
    (\vec{a} + \vec{b})^T (\vec{a} - \vec{b}) &= (\vec{a}^T + \vec{b}^T)(\vec{a} - \vec{b}) \nonumber \\
    &= \vec{a}^T\vec{a} - \vec{a}^T\vec{b} + \vec{b}^T\vec{a} - \vec{b}^T\vec{b}
\end{align}
Since the dot product is commutative, $\vec{a}^T\vec{b} = \vec{b}^T\vec{a}$. The expression simplifies to:
\begin{align}
    (\vec{a} + \vec{b})^T (\vec{a} - \vec{b}) = \vec{a}^T\vec{a} - \vec{b}^T\vec{b} = |\vec{a}|^2 - |\vec{b}|^2 \label{eq:algebra}
\end{align}
Now, we calculate the squared magnitudes of $\vec{a}$ and $\vec{b}$ using their numerical values.
The given vectors are:
\begin{equation}
\vec{a} = \myvec{5 \\ -1 \\ -3}, \quad \vec{b} = \myvec{1 \\ 3 \\ -5}
\end{equation}
The squared magnitudes are:
\begin{align}
    |\vec{a}|^2 = \vec{a}^T\vec{a} &= (5)^2 + (-1)^2 + (-3)^2 \nonumber \\
    &= 25 + 1 + 9 = 35 \label{eq:mag_a} \\
    |\vec{b}|^2 = \vec{b}^T\vec{b} &= (1)^2 + (3)^2 + (-5)^2 \nonumber \\
    &= 1 + 9 + 25 = 35 \label{eq:mag_b}
\end{align}
Substituting the results from \eqref{eq:mag_a} and \eqref{eq:mag_b} into our algebraic result from \eqref{eq:algebra}:
\begin{align}
    (\vec{a} + \vec{b})^T (\vec{a} - \vec{b}) = 35 - 35 = 0
\end{align}
Since the dot product is zero, the vectors $(\vec{a} + \vec{b})$ and $(\vec{a} - \vec{b})$ are perpendicular. Hence, proved.

\begin{figure}[H]
    \centering
    \includegraphics[width=1.1\columnwidth]{figs/Figure_1.png}
    \label{fig:placeholder}
\end{figure}

\end{document}